\documentclass[12pt, letterpaper]{article}

%-------Packages---------
\usepackage{amssymb,amsfonts}
\usepackage{mathptmx}
\usepackage{amsthm}
\usepackage{graphicx}
\usepackage[all,arc]{xy}
\usepackage{enumerate}
\usepackage{bbm}
\usepackage{dsfont}
\usepackage{mathrsfs}
\usepackage{tikz-cd}
\usepackage{setspace}
\usepackage{import}
\usepackage[utf8]{inputenc}
\usepackage{hyperref}
\usepackage{tikz}
\usepackage{float}
\usepackage[dvipsnames]{xcolor}
\usepackage{fancyhdr}
\usepackage[nointegrals]{wasysym}

\usepackage[left=1in,top=1in,right=1in,bottom=1in]{geometry}

\usepackage{mathtools}
\usepackage{setspace}
\usepackage{cleveref}
\usepackage{multicol}
\usepackage{mathpazo}
\usepackage{tcolorbox}
\tcbuselibrary{theorems,skins,breakable}
\usepackage{xparse}

\usepackage{xcolor, ifthen, tcolorbox}
\tcbuselibrary{theorems,skins,breakable}
% Theorem System
% The following environments are provided:
%   New Problem:        \begin{problem} ...
%   New Question Header:   \qh (then followed by subproblems)
%   Subproblem:     \begin{subproblem} ...
%   Problem Proof:  \begin{problemproof} ...
%   Proof:          \\begin{pf}{color} ...
%   Remark:         \rmk (sentence), \rmkb (block)

% Define custom colors
\definecolor{thmscol}{HTML}{F4565B} %For Problems
\definecolor{lemscol}{HTML}{FE844C} %For Lemmes
\definecolor{prpscol}{HTML}{00A7EF} %For proposition
\definecolor{corscol}{HTML}{23DAFF} %For corrolaries
\definecolor{rmkscol}{HTML}{6DEEC5} %For remarks
\definecolor{defscol}{HTML}{18C991} %For definitions
\definecolor{exmscol}{HTML}{7DB800} %For examples
\definecolor{clmscol}{HTML}{165c58} %For claims

%Change between dark(black)/light(white) mode
\newcommand{\colormode}{white}
\ifthenelse{\equal{\colormode}{white}}
      {\newcommand{\TextColor}{black}}
      {\newcommand{\TextColor}{white}}
\newcommand{\pbcolormain}{thmscol!80!\colormode}
\newcommand{\pbcolorsecond}{thmscol!38!\colormode}


% ============================
% Problem Counter
% ============================
\newcounter{subproblemcounter}
\newcounter{problemcounter}

\newcommand{\q}{
    \setcounter{subproblemcounter}{0}%
    \stepcounter{problemcounter} % Increment problem counter
    \color{\TextColor}Problem \arabic{problemcounter}: % Display the problem counter
}

\newcommand{\stepsub}{
    \stepcounter{subproblemcounter}%
    \color{\TextColor}(\alph{subproblemcounter})
}
% ============================
% Problem
% ============================
\newtcolorbox{pb}[1]{
  enhanced,
  frame hidden,
  titlerule=0mm,
  toptitle=1mm,
  bottomtitle=1mm,
  fonttitle=\bfseries\large,
  coltitle=black,
  colbacktitle=\pbcolormain,
  colback=\pbcolorsecond,
  title={\q #1},
  coltext=\TextColor
}

\newtcolorbox{spb}[1]{
  enhanced,
  frame hidden,
  titlerule=0mm,
  toptitle=1mm,
  bottomtitle=1mm,
  fonttitle=\bfseries\large,
  coltitle=black,
  colbacktitle=\pbcolormain,
  colback=\pbcolorsecond,
  title={\stepsub #1},
  coltext=\TextColor,
  attach boxed title to top left={xshift=3mm,yshift=-2mm},
  boxed title style={
    colback=\pbcolormain,
    colframe=\pbcolormain,
  }
}

\newtcolorbox{pbh}[1]{
    enhanced,
    frame hidden,
    titlerule=0mm,
    toptitle=1mm,
    bottomtitle=1mm,
    fonttitle=\bfseries\large,
    coltitle=black,
    colbacktitle=\pbcolormain,
    colback=\pbcolorsecond,
    title={\q #1},
    boxrule=0pt,
    interior hidden,
    attach boxed title to top left={xshift=3mm,yshift=-2mm},
    boxed title style={
        colback=\pbcolormain,
        colframe=\pbcolormain,
    }
}

\newenvironment{problem}[1]{%
  \newpage
  \begin{pb}{#1}
    }{
  \end{pb}
}

\newenvironment{subproblem}[1]{%
  \begin{spb}{#1}
    }{
  \end{spb}
}

\newenvironment{problemproof}{%
    \begin{pf}{\pbcolormain}
        }{
    \end{pf}
}

\newcommand{\qh}[1][]{
    \newpage
    \begin{pbh}{#1}\end{pbh} % Display the problem counter
}

% ============================
% Claim
% ============================

\newtcbtheorem[number within=problemcounter]{claim}{Claim}
{
    enhanced,
    frame hidden,
    titlerule=0mm,
    toptitle=1mm,
    bottomtitle=1mm,
    fonttitle=\bfseries\large,
    coltitle=black,
    colbacktitle=clmscol!50!\colormode,
    colback=clmscol!20!\colormode,
}{clm}

\NewDocumentCommand{\clm}{m+m}{
    \begin{claim*}{#1}{}
        #2
    \end{claim*}
}

\newenvironment{clmpf}{
	{\noindent{\it \textbf{Proof.}}
	\setlength{\parindent}{0pt}
	}
	\tcolorbox[blanker,breakable,left=5mm,parbox=false,
    before upper={\parindent15pt},
    after skip=10pt,
	borderline west={1mm}{0pt}{clmscol!50!\colormode}]
}{
    \textcolor{clmscol!50!\colormode}{\hbox{}\nobreak\hfill$\blacksquare$}
    \endtcolorbox
}

\NewDocumentCommand{\clmp}{m+m+m}{
    \begin{claim*}{#1}{}
        #2
    \end{claim*}

    \begin{clmpf}
        #3
    \end{clmpf}
}


% ============================
% Proof
% ============================

\newenvironment{pf}[1]{%
  \def\pfcolor{#1}% Store the color in a macro
  \noindent{\it \textbf{Proof.}}%
  \tcolorbox[blanker,breakable,left=5mm,parbox=false,%
    before upper={\parindent15pt},%
    after skip=10pt,%
    borderline west={1mm}{0pt}{\pfcolor},%
    coltext=\TextColor
  ]%
  \setlength{\parindent}{0pt}%
}{%
  \textcolor{\pfcolor}{\hbox{}\nobreak\hfill$\blacksquare$}%
  \endtcolorbox%
}

\newtcolorbox{solution}{
  blanker,
  breakable,
  left=5mm,
  parbox=false,%
  before upper={\parindent15pt},%
  after skip=10pt,%
  borderline west={1mm}{0pt}{\pbcolormain},%
  coltext=\TextColor
}


% ============================
% Remark
% ============================


\NewDocumentCommand{\rmk}{+m}{
    {\it \color{rmkscol!100!white}#1}
}

\newenvironment{remark}{
    \par
    \vspace{5pt}
    \begin{minipage}{\textwidth}
        {\par\noindent{\textbf{Remark.}}}
        \tcolorbox[blanker,breakable,left=5mm,
        before skip=10pt,after skip=10pt,
        borderline west={1mm}{0pt}{rmkscol!80!\colormode},
        coltext=\TextColor
        ]
}{
        \endtcolorbox
    \end{minipage}
    \vspace{5pt}
}

\NewDocumentCommand{\rmkb}{+m}{
    \begin{remark}
        #1
    \end{remark}
}


% Define a new command for problems
\renewcommand{\headrule}{\color{\TextColor}\hrulefill}
\newcommand{\C}{\mathbb{C}}
\newcommand{\R}{\mathbb{R}}
\newcommand{\Q}{\mathbb{Q}}
\newcommand{\Z}{\mathbb{Z}}
\newcommand{\N}{\mathbb{N}}
\renewcommand{\P}{\mathbb{P}}
\newcommand{\F}{\mathbb{F}}
\newcommand{\T}{\mathcal{T}}
\newcommand{\contr}{\Rightarrow \Leftarrow}
\newcommand{\s}{\(\,\)}
\renewcommand{\t}[1]{\text{#1}}
\newcommand{\ang}[1]{\left\langle #1 \right\rangle}
\newcommand{\coker}{\mathrm{coker}}

% Meta data
\setstretch{1.2}

\begin{document}
\color{\TextColor}
\pagecolor{\colormode}
\setlength{\parindent}{0pt}
\pagestyle{fancy}
\fancyhf{}
\fancyhead[LOE]{\color{\TextColor}\slshape{\textbf{Vincent Ferrara}}}
\fancyhead[ROE]{\color{\TextColor}\thepage}

\begin{problem}{}
    Let $X$ be a regular space. Show that for every pair of distinct
points $x, y$ in $X$, there are neighborhoods $U_x, U_y$ of $x, y$
with $\overline{U_x} \cap \overline{U_y} = \emptyset$.
\end{problem}
\begin{problemproof}
    Let $x\neq y \in X$. Since $X$ is regular this means $\{y\}$ is closed.

    Thus, there is disjoint open sets s.t. $x \in U$ and $\{y\} \subset V$.

    Thus, proven in class since $X$ is regular there exists neighborhoods $U_x$, $U_y$ s.t. $\overline{U_x} \subset U$ and $\overline{U_y} \subset V$.

    Since $U \cap V = \emptyset \implies \overline{U_x} \cap \overline{U_y} = \emptyset$.
\end{problemproof}

\begin{problem}{}
    Let $\{X_\alpha\}_{\alpha \in J}$ be a collection of nonempty
spaces, and suppose that the product $\prod_{\alpha \in J}X_\alpha$
is normal. Show that $X_\alpha$ is normal for each $\alpha \in
J$. (Note: the converse to this statement is false, but a
counterexample is difficult to come by; see Exercise 9 in \S 32 of
Munkres.)
\end{problem}
\clmp{}{
    Let $X=\prod_{\alpha \in J} X_\alpha$ be a product of topological spaces with the product topology.

    Let $\pi_i : X \to X_i$ be a projection map.

    $\pi_i$ is an open map.
}{
    Let $B=\prod U_\alpha$ be a basic open set in $X$ s.t. $U_\alpha$ is open in $X_\alpha$ and $U_\alpha = X_\alpha$ for all but finitely many indices.

    \noindent Consider $\pi_i(B)$.

    \noindent If $i$ is an index where $U_i \neq X_i$ then $\pi_i(B)=U_i$ which is open

    \noindent If $i$ is an index where $U_i = X_i$ then $\pi_i(B)=X_i$ which is open.

    \noindent Thus, $\pi_i$ is open on the basis sets.

    \noindent Let $U$ be open in $X$. Thus, $U=\bigcup B_\beta$ s.t. $B_\beta$ is a basis element. Thus,
    \[\pi_i(U)=\pi_i(\bigcup B_\beta)=\bigcup \pi_i(B_\beta)\]
    Since $\pi_i(B\beta)$ is open in $X_i$, it follows that the union of all of these which is $\pi_i(U)$ is open.

    \noindent Thus, projection maps are open.
}
\begin{problemproof}
    Let $X = \prod_{\alpha \in J} X_\alpha$

    Let $X_\beta$ be one element in the collection of spaces.

    Let $\pi_\beta$ be the projection map from the product to this one space.

    Let $A,B \subset X_\alpha$ be disjoint closed subsets.

    Since $\pi_\beta$ is continuous we know that $\pi^{-1}_\beta(A),\pi^{-1}_\beta(B)$ are closed in $X$.

    Thus, since $A \cap B = \emptyset$. This implies that $\pi^{-1}_\beta(A) \cap \pi^{-1}_\beta(B) = \emptyset$. Thus, since $X$ is normal there exists disjoint open subsets $W_A$, $W_B$ in $X$ s.t. $\pi^{-1}_\beta(A) \subset W_A$ and $\pi^{-1}_\beta(B) \subset W_B$.

    Thus, by the proof above we get that $U=\pi_\beta(W_A)$ and $V=\pi_\beta(W_B)$ are disjoint open sets in $X_\beta$ that contain $A$, $B$ respectively.

    Also, since $X$ is normal this implies it is Hausdorff. Thus, proven in class we have shown that each component is thus Hausdorff, thus each component is also $T_1$.
    
    Thus, each component is normal.
\end{problemproof}

\begin{problem}{}
    Show that a connected normal space containing at least two points
contains an uncountable amount of points.
\end{problem}
\begin{problemproof}
    Let $X$ be a connected normal space with points $x\neq y \in X$.

    By Urysohn Lemma we know there exists a continuous function $f: X \to [0,1]$ s.t. $f(x)=0$ and $f(y)=1$. We use the closed subsets $\{x\}$ and $\{y\}$ since $X$ is also $T_1$.

    Thus, since $X$ is connected we know that $f(X) \subset [0,1]$ is connected. Thus, by the intermediate value theorem we know since $0,1 \in f(X)$ then $[0,1] \subset f(X)$. Thus, $f(X)=[0,1]$.

    Therefore, $X$ is uncountable since $f$ surjects $X$ onto an uncountable space.
\end{problemproof}

\begin{problem}{}
    A topological space $X$ is said to be \emph{completely regular} if
finite subsets of $X$ are closed, and $X$ satisfies the following
condition: for every point $x \in X$ and every closed subset
$A \subset X$, there exists a continuous function $f:X \to [0,1]$
such that $f(x) = 0$ and $f(a) = 1$ for every $a \in A$.
\end{problem}

\begin{subproblem}{}
    Explain why a completely regular space is always regular, and
why a normal space is always completely regular
\end{subproblem}
\begin{problemproof}
    Let $X$ be completely regular. Let $x \in X$ and let $A \subset X$ be a closed set s.t. $x \notin A$.

    Then, there exists a continuous function $f: X \to [0,1]$ s.t. $f(x)=0$ and $f(a)=1$ for all $a \in A$.

    Consider $x \in U=f^{-1}([0,1/2))$ and $A \subset V=f^{-1}((1/2,1])$ these are open sets.

    Since $[0,1/2)$ is disjoint from $(1/2,1]$, This implies that $U \cap V = \emptyset$. Thus, $X$ is regular.

    Let $X$ be normal. Then by Urysohn Lemma we know that there exists a continuous function $f: X \to [0,1]$ s.t. $f(x)=0$ and $f(a)=1$ for all $a \in A$.

    Thus, $X$ is completely regular.
\end{problemproof}

\begin{subproblem}{}
    Prove that every locally compact Hausdorff space is
completely regular. (Use the one-point compactification.)
\end{subproblem}
\begin{problemproof}
    Let $X$ be a locally compact Hausdorff space. Thus, we know there exists a one point compactification $\hat{X}=X \cup \{\infty\}$.

    Thus, $\hat{X}$ is a compact Hausdorff spaces proven in class we know that $\hat{X}$ is normal.

    Let $x \in X$ and $A \subset X$ be a closed subset s.t. $x \notin A$. Thus, $X \setminus A$ is open in $X$ and $\hat{X}$. Therefore, $\hat{X} \setminus (X \setminus A)=A \cup \{\infty\}$ is closed in $\hat{X}$.

    Since $\hat{X}$ is normal by the problem above we get that it is completely regular thus there is a function $f: \hat{X} \to [0,1]$ s.t. $f(x)=0$ and $f(a)=1$ for all $a \in A \cup \{\infty\}$.

    Now consider $p: X \to \hat{X}$ s.t. $p(x)=x$. Let $O$ be open in $\hat{X}$ then if $O$ is already open in $X$ then $p^{-1}(O)=O$.
    
    If $O=\hat{X} \setminus K$ for some compact set $K$ in $X$ then $p^{-1}(O)=X \setminus K$. Since $X$ is Hausdorff and $K$ is compact this means that $K$ is closed thus $X \setminus K$ is open.

    Thus, $p$ is continuous.

    Thus, consider $f \circ p: X \to [0,1]$. This is a continuous function s.t. $f(p(x))=f(x)=0$, and $f(p(a))=f(a)=1$ for all $a \in A \subset A \cup \{\infty\}$.

    Thus, $X$ is completely regular.
\end{problemproof}

\begin{problem}{Topological groups}
    Recall that a group $G$ is just a set equipped with some extra
structure: a binary operation satisfying certain axioms. A
topological space $X$ is also a set equipped with some extra
structure: a topology $\T$, or a collection of subsets of $X$.
A \emph{topological group} is a set which is equipped with the
structure of \emph{both} a group and a topological space, in a
``compatible'' way. Precisely, a topological group is a group $G$
which is also a topological space such that:
\begin{enumerate}[I]
    \item The maps
    \[
    G \times G \to G, \qquad G \to G
    \]
    given respectively by $(g, h) \mapsto gh$ and $g \mapsto g^{-1}$
    are both continuous, and
    \item The singleton $\{e\}$ is a closed subset of $G$, where $e$
    is the identity element.
\end{enumerate}
For the below, let $G$ be a topological group.
\end{problem}

\begin{subproblem}{}
    Prove that any topological group satisfies the $T_1$ axiom.
\end{subproblem}
\clmp{}{
    Let $X,Y$ be topological spaces then $f: X \to X \times \{y\}$ for $y \in Y$ is a homeomorphism.
}{
    Injective:

    \noindent Let $a,b \in X$ s.t. $f(a)=f(b) \implies (a,y)=(b,y) \implies a=b$

    \noindent Surjective:

    \noindent Let $(a,y) \in X \times \{y\}$. Consider $f(a)=(a,y)$

    \noindent Continuous:

    \noindent Let $U$ be a basic open set in $X \times \{y\}$. Thus, $U=V \times W$ for $V$ open in $X$ and $W$ open in $\{y\}$.

    \noindent Thus, $f^{-1}(U)=V$.

    \noindent Continuous Inverse:

    \noindent Let $U$ be an open set in $X$. Thus, $f(U)=U \times \{y\}$. This is a basic open set in $X \times \{y\}$. Thus, $f(U)$ is open.

    \noindent Therefore, $f$ is a homeomorphism.
}
\begin{problemproof}
    Let $m : G \times G \to G$ be the binary operation map.

    Define: $f_x : G \to \{x\} \times G$ s.t. $f_x(g)=(x,g)$

    Define: $L_x : G \to G$ s.t. $L_x = m \circ f_x$

    Injective:

    Let $a,b \in G$ s.t. $L_x(a)=L_x(b) \implies xa=xb$. Multiply by $x^{-1}$ on both sides, and we get that $a=b$.

    Surjective:

    Let $a \in G$ consider $y=x^{-1}a$

    Thus, $L_x(y)=xx^{-1}a=a$

    Continuous:

    $L_x$ is continuous since $f_x$ by the claim and proof above and $m$ is continuous by definition.

    Continuous Inverse:

    Since the inverse of $L_x^{-1}=L_{x^{-1}}$ is also continuous by the same proof above we get that $L_x$ is a homeomorphism.

    Therefore, we get that $L_x(\{e\})=\{x\}$ is closed.
\end{problemproof}

\begin{subproblem}{}
    Prove that any topological group is Hausdorff.
\end{subproblem}
\begin{problemproof}
    Let $i : G \to G$ be the inverse map.

    Define: $f : G \times G \to G$ s.t. $f(g,h)=m(g,i(h))=gh^{-1}$. This is continuous since $m,i$ and the identity map are continuous.

    Now consider $f^{-1}(\{e\})=\{(g,h) \in G \times G \mid gh^{-1}=e\}$. Since invserse are unique in groups this is only true when $g=h$.

    Thus, we have shown that $\Delta=\{(g,g) \mid g \in G\}$ is closed. Thus, by previous homework $G$ is Hausdorff.
\end{problemproof}

\begin{subproblem}{}
    Prove that any topological group is regular.
\end{subproblem}
\begin{problemproof}
    Consider the continuous map from the last problem $f(g,h)=gh^{-1}$.

    Let $A$ be a closed set s.t. $e \notin A$.

    Consider $f^{-1}(X \setminus A)$. This is an open set in $G \times G$ since $f$ is continuous.

    Since $f(e,e)=e \in X \setminus A \implies (e,e) \in f^{-1}(X \setminus A)$. Therefore, let $V_1 \times V_2$ be a basic open set s.t. $(e,e) \in V_1 \times V_2 \subset f^{-1}(X \setminus A)$.

    Define: $V=V_1 \cap V_2$. This is open since $V_1,V_2$ are open in $G$.

    Consider $AV=\displaystyle\bigcup_{a \in A}L_a(V)$. This is an open set since $L_a(V)$ is open since $L_a$ is an homeomorphism.

    Let $a \in A$ since $e \in V \implies ae=a \in AV$. Thus, $A \subset V$.

    Let $x \in V$ and $x \in AV$. Thus, $x=ay$ for some $a \in A$ and $y \in Y$. Thus, $a=xy^{-1}$. This is a contradiction since $x,y \in V \implies f(x,y) \in X \setminus A \implies a \in X \setminus A$.

    Therefore, $V \cap AV = \emptyset$, and $e \in V$, $A \subset AV$.

    To do this for any $x \in G$ and any closed set $C$ s.t. $x \notin C$. Just do this for $e$, and $x^{-1}C$. $e \notin x^{-1}C$ since if it was then this would imply that $x \in C$ since inverses are unique.

    Thus, we get disjoint open sets s.t. $e \in V$ and $x^{-1}C \subset U$.

    Then consider $x \in xV$ and $C \subset xU$. These are still disjoint since if there were not this would imply that $V \cap U \neq \emptyset$ since left multiplication is a homeomorphism.

    Thus, $G$ is regular.
\end{problemproof}
\end{document}